\chapter{Glossary}
\label{app:glossary}

\begin{description}[leftmargin=0em, style=nextline]
  \item[active space] The orbital window inside which electron correlation is
    treated explicitly (CASSCF's ``CAS''). The program suggests starting
    windows and checks the occupations; the choice stays yours.
  \item[CASSCF] Complete-active-space self-consistent field: orbitals are
    optimised together with a full CI inside the active space. Its convergence
    has two different criteria (energy and orbital gradient) -- the difference
    matters for everything computed on top.
  \item[NEVPT2 / CASPT2] Second-order perturbation corrections on top of a
    CASSCF reference (the ``dynamic correlation''). They depend on converged
    \emph{orbitals}, and each has its own failure signature (intruder states).
  \item[SC-NEVPT2] Strongly contracted NEVPT2 -- the variant generated in this
    project's ORCA inputs.
  \item[ECP / pseudopotential] An effective core potential replacing the core
    electrons; for lanthanides/actinides it also carries scalar-relativistic
    effects. 5f-in-core variants remove the f electrons too -- legitimate only
    when no f physics is wanted.
  \item[single-orbital entropy] An entanglement measure of one orbital with all
    the others, defined on the four occupation states of that orbital
    (ceiling $\ln 4$). An ORCA output does not determine it; the program's A2
    section computes a strict \emph{upper bound} and uses it one-sidedly.
  \item[T1 diagnostic] The norm of the coupled-cluster singles amplitudes
    (scaled by the correlated electron count); above about 0.02 for
    closed-shell references it flags that a single-reference treatment is
    questionable. ORCA prints T1, not D1.
  \item[crystal-field parameters $B_k^q$] The coefficients of the
    crystal-field Hamiltonian in the $|J M\rangle$ manifold, in an operator
    convention that must be declared (Stevens lineage in this program).
    Values from different conventions, manifolds, units or $z$-axes are not
    comparable parameter by parameter.
  \item[Stevens operators] Extended operator equivalents $O_k^q(J)$ -- the
    operator set the crystal-field fit uses (transcribed from the EasySpin
    documentation table; see the module's source list).
  \item[l-shell vs $J$-manifold parameters] The point-charge estimate
    (menu 11) produces exact $l = 3$ matrix elements ($A_k^q \langle
    r^k\rangle$); the fit (menu 10) produces $J$-manifold Stevens parameters.
    The projection between them is not performed here; compare spectra, not
    parameters.
  \item[$\Delta S_E$] The environment spin-polarisation entropy: how far the
    environment density is from the spin-unpolarised point. The real quantity
    needs spin-resolved 2-RDMs (\code{orca\_2json}) plus localised orbitals
    (\code{orca\_loc}); the program's A6 section reports ORCA's local spin
    analysis as a clearly-labelled surrogate.
  \item[local spin analysis] ORCA's fragment-resolved spin expectation values
    ($\langle S_z^A\rangle$, $S_{\mathrm{eff}}(A)$, $\langle S_A S_B\rangle$),
    printed when the input divides the molecule into fragments.
  \item[fixture] A registered real output file against which the parsers and
    reports are tested; provenance and generation are recorded next to it.
  \item[provenance classes] manual (section + URL), literature (DOI-backed
    citation + BibTeX), measured (a record of this project, by path).
  \item[refusal] A first-class outcome: the program declining to conclude, with
    the reason and the next step. Refusals are results, not failures.
\end{description}
